\section{Objetivos}

\subsection{Objetivo general}

Analizar la relación entre la presencia de micronegocios y el nivel de empleo a nivel municipal en el estado de Yucatán, utilizando herramientas de econometría y análisis espacial, con el fin de evaluar si estos establecimientos actúan como motor del empleo o responden a condiciones estructurales del mercado laboral.

\subsection{Objetivos específicos}

\begin{itemize}

\item Construir una base de datos a nivel municipal a partir de los Censos Económicos 2019 del \cite{inegi2019} y el Directorio Estadístico Nacional de Unidades Económicas (DENUE) \cite{denue}, que permita identificar el número de micronegocios y el empleo total.

\item Analizar la distribución espacial de los micronegocios y el empleo en Yucatán mediante herramientas de sistemas de información geográfica (QGIS), incluyendo mapas coropléticos y de densidad.

\item Identificar patrones de concentración geográfica en la actividad económica mediante técnicas de análisis espacial, como la autocorrelación espacial.

\item Estimar un modelo econométrico a nivel municipal que permita evaluar la relación entre el número de micronegocios y el empleo total.

\item Interpretar los resultados obtenidos para determinar si los micronegocios tienen un efecto significativo en la generación de empleo o si su presencia está asociada a otras condiciones estructurales.

\end{itemize}